\chapter{Appendix: Sets, Counting, and Probability}

\providecommand{\E}{\ensuremath{\mathbb{E}}}
\providecommand{\Var}{\ensuremath{\operatorname{Var}}}

This is a toolbox chapter. The rest of the book leans on it constantly --- every
time we say ``there are $n!$ orderings,'' ``the expected running time is $O(n\log n)$,''
or ``the probability of a bad pivot is small,'' we are quietly using the language
built here. CLRS folds this material into its appendices and expects you to already
be comfortable with it. We are not going to assume that. We are going to build it
slowly, in plain words, and then make it precise.

You can read this chapter front to back, or you can keep it open beside you as a
reference and dip in when a later chapter says ``recall expectation'' or ``by
linearity.'' Either way, the goal is the same: by the end you should be able to
\emph{count} the things an algorithm does and \emph{reason about chance} when an
algorithm flips coins. Those two skills --- counting and probability --- are the
mathematical heart of algorithm analysis.

\section{Sets: the language everything is written in}
\label{ch38:sets}

\begin{intuitionbox}
A \emph{set} is just a collection of distinct things, with no order and no
repetition. Think of it as a bag where you only care \emph{which} items are inside,
not how many copies there are or what order you pulled them out. The bag
$\{3, 1, 2\}$ is the very same bag as $\{1, 2, 3\}$ and as $\{1, 1, 2, 3\}$ --- the
duplicate and the reordering simply do not register.
\end{intuitionbox}

We write $x \in A$ to say ``$x$ is an element of $A$'' and $x \notin A$ for the
negation. The set with nothing in it is the \emph{empty set}, written $\emptyset$ or
$\{\}$. Some sets you already know:
\[
\N = \{0, 1, 2, 3, \dots\}, \qquad
\Z = \{\dots, -2, -1, 0, 1, 2, \dots\}, \qquad
\R = \text{the real numbers.}
\]
(We follow the convention that $\N$ includes $0$; some authors start at $1$. When it
matters, we will say so.)

A set can be described by \emph{listing} its elements, $\{2, 4, 6, 8\}$, or by giving
a \emph{rule}, the set-builder notation
\[
\{\, x \in \N : x \text{ is even and } x \le 8 \,\}.
\]
Read the colon as ``such that.'' Both descriptions name the same four-element set.

\subsection{Size, subsets, and the empty set}

The number of elements in a finite set $A$ is its \emph{cardinality}, written
$\abs{A}$. So $\abs{\{2,4,6,8\}} = 4$ and $\abs{\emptyset} = 0$. Cardinality is the
single most important thing this chapter teaches you to compute, because the running
time of an algorithm is almost always ``do a fixed amount of work for each element of
some set,'' and so the time is ``$\abs{\text{that set}}$ times a constant.''

We say $A$ is a \emph{subset} of $B$, written $A \subseteq B$, when every element of
$A$ is also an element of $B$. The empty set is a subset of \emph{every} set (there
is no element of $\emptyset$ to check, so the condition holds vacuously), and every
set is a subset of itself. If $A \subseteq B$ but $A \ne B$ we call $A$ a
\emph{proper} subset and write $A \subset B$.

\begin{trapbox}
Do not confuse $\in$ with $\subseteq$. The number $3$ is an \emph{element} of
$\{1,2,3\}$, so $3 \in \{1,2,3\}$. The one-element set $\{3\}$ is a \emph{subset},
so $\{3\} \subseteq \{1,2,3\}$. Writing $3 \subseteq \{1,2,3\}$ or
$\{3\} \in \{1,2,3\}$ is a type error --- the same kind of mistake as adding a
pointer to an integer. Keep ``is one of the things inside'' ($\in$) separate from
``is a smaller bag inside'' ($\subseteq$).
\end{trapbox}

\subsection{Operations on sets}

Given two sets $A$ and $B$, we combine them in four standard ways. Picture two
overlapping circles --- a Venn diagram --- as you read each one.

\begin{itemize}
\item \textbf{Union} $A \cup B$: everything in $A$, or in $B$, or in both.
\item \textbf{Intersection} $A \cap B$: only the things in \emph{both}.
\item \textbf{Difference} $A \setminus B$: the things in $A$ that are \emph{not} in $B$.
\item \textbf{Complement} $\overline{A}$: everything (in some agreed-upon universe $U$)
that is \emph{not} in $A$. So $\overline{A} = U \setminus A$.
\end{itemize}

\begin{center}
\begin{tikzpicture}[scale=1.0]
  \begin{scope}
    \clip (-0.6,-1.2) rectangle (3.0,1.2);
    \fill[sourceblue!18] (0,0) circle (1.1);
    \fill[targetgreen!18] (1.4,0) circle (1.1);
    \begin{scope}
      \clip (0,0) circle (1.1);
      \fill[structurepurple!28] (1.4,0) circle (1.1);
    \end{scope}
  \end{scope}
  \draw[thick,sourceblue] (0,0) circle (1.1);
  \draw[thick,targetgreen] (1.4,0) circle (1.1);
  \node[sourceblue!60!black] at (-0.35,0) {$A$};
  \node[targetgreen!60!black] at (1.75,0) {$B$};
  \node[structurepurple!70!black] at (0.7,0) {\scriptsize $A\cap B$};
  \draw (-0.6,-1.2) rectangle (3.0,1.2);
  \node at (2.7,1.0) {$U$};
\end{tikzpicture}
\end{center}

Two sets with empty intersection, $A \cap B = \emptyset$, are called
\emph{disjoint}: their circles do not overlap. Disjointness is the key that unlocks
clean counting, as we are about to see.

\begin{formalbox}
\textbf{De Morgan's laws.} For any sets in a universe $U$,
\[
\overline{A \cup B} = \overline{A} \cap \overline{B},
\qquad
\overline{A \cap B} = \overline{A} \cup \overline{B}.
\]
In words: ``not (in $A$ or $B$)'' means ``not in $A$ and not in $B$,'' and ``not (in
both)'' means ``missing from at least one.'' These mirror the logic laws
$\lnot(p \lor q) = \lnot p \land \lnot q$ exactly, because a set is just a predicate
in disguise.
\end{formalbox}

\subsection{Ordered pairs, tuples, and the Cartesian product}

Sets forget order. But algorithms care about order all the time --- a coordinate
$(x, y)$ is not the same point as $(y, x)$. So we also need \emph{ordered} structure.

An \emph{ordered pair} $(a, b)$ remembers which came first. Two pairs are equal only
when both slots match: $(a,b) = (c,d)$ exactly when $a = c$ and $b = d$. More
generally an \emph{$n$-tuple} $(a_1, a_2, \dots, a_n)$ is an ordered list of $n$
slots.

The \emph{Cartesian product} $A \times B$ is the set of \emph{all} ordered pairs with
first slot from $A$ and second from $B$:
\[
A \times B = \{\, (a, b) : a \in A,\ b \in B \,\}.
\]

\begin{example}[A small product]
\label{ch38:ex-product}
Let $A = \{1, 2\}$ and $B = \{x, y, z\}$. Then
\[
A \times B = \{(1,x),(1,y),(1,z),(2,x),(2,y),(2,z)\},
\]
which has $2 \times 3 = 6$ elements. The grid you would draw to list them --- $2$
rows, $3$ columns --- is exactly why $\abs{A \times B} = \abs{A}\cdot\abs{B}$. Hold
onto that picture; it is the seed of the entire counting section.

\begin{center}
\begin{tikzpicture}[scale=1.0,every node/.style={font=\small}]
  \node[sourceblue!60!black] at (-0.8,0.6) {$A$};
  \node[targetgreen!60!black] at (0.6,1.7) {$B$};
  % column headers (B)
  \foreach \c/\b in {0/x, 1/y, 2/z} {
    \node[targetgreen!60!black] at (\c,1.2) {$\b$};
  }
  % row headers (A)
  \foreach \r/\a in {0/2, 1/1} {
    \node[sourceblue!60!black] at (-0.8,{1-\r}) {$\a$};
  }
  % the 6 product cells
  \foreach \r/\a in {0/2, 1/1} {
    \foreach \c/\b in {0/x, 1/y, 2/z} {
      \node[draw,structurepurple!60!black,fill=structurepurple!8,
            rounded corners,minimum width=1.0cm,minimum height=0.7cm]
        at (\c,{1-\r}) {\scriptsize $(\a,\b)$};
    }
  }
\end{tikzpicture}
\end{center}
The $\abs{A} = 2$ rows times the $\abs{B} = 3$ columns give one cell per pair: a
$2 \times 3 = 6$ rectangle of outcomes.
\end{example}

\section{Relations and functions}
\label{ch38:relations}

\subsection{Relations}

A \emph{relation} from $A$ to $B$ is any subset $R \subseteq A \times B$. That sounds
abstract, so unpack it: a relation is a rule that, for each pair $(a,b)$, says either
``yes, these are related'' (the pair is in $R$) or ``no'' (it is not). When
$(a,b) \in R$ we often write $a\, R\, b$. The familiar $\le$ on numbers is a
relation: the pair $(3, 5)$ is in it because $3 \le 5$; the pair $(5, 3)$ is not.

A relation $R$ on a single set $A$ (meaning $R \subseteq A \times A$) can have nice
properties:
\begin{itemize}
\item \textbf{Reflexive}: $a\, R\, a$ for every $a$. (Everything relates to itself.)
\item \textbf{Symmetric}: whenever $a\, R\, b$, also $b\, R\, a$.
\item \textbf{Transitive}: whenever $a\, R\, b$ and $b\, R\, c$, also $a\, R\, c$.
\end{itemize}
A relation that is reflexive, symmetric, and transitive all at once is an
\emph{equivalence relation}. It carves $A$ into disjoint clumps --- \emph{equivalence
classes} --- where everything in a clump is related to everything else in that clump
and to nothing outside it. ``Has the same remainder mod $7$'' is an equivalence
relation on $\Z$; it sorts the integers into $7$ classes. Equivalence classes are the
mathematical backbone of union-find (Chapter on disjoint sets), where the whole job
is to track which clump each element belongs to.

\subsection{Functions}

\begin{intuitionbox}
A \emph{function} is a machine: feed it one input, it gives you back exactly one
output. The crucial word is \emph{exactly one}. A vending machine that sometimes
returns chips and sometimes a soda for the same button-press is not a function; a
machine that returns nothing for some buttons is not (a total) function either. One
input in, one output out, every time.
\end{intuitionbox}

Formally a function $f : A \to B$ is a special relation: for each $a \in A$ there is
exactly one $b \in B$ with $(a, b) \in f$, and we write $b = f(a)$. The set $A$ is the
\emph{domain} (allowed inputs), $B$ is the \emph{codomain} (the type of the outputs),
and the \emph{image} (or range) is the set of values actually hit,
$\{\, f(a) : a \in A \,\}$.

Three properties matter constantly in algorithms:
\begin{itemize}
\item \textbf{Injective} (one-to-one): different inputs give different outputs. If
$f(a_1) = f(a_2)$ then $a_1 = a_2$. No collisions.
\item \textbf{Surjective} (onto): every element of the codomain is hit by something.
The image equals the whole codomain.
\item \textbf{Bijective}: both injective and surjective. A perfect pairing --- every
input matched to a unique output and vice versa.
\end{itemize}

\begin{focusbox}
Bijections are the secret weapon of counting. If you can build a bijection between
two finite sets, you have \emph{proved} they are the same size without counting
either one. ``Show $\abs{A} = \abs{B}$'' becomes ``pair them up perfectly.'' We will
use this trick in the next section, and it appears again whenever a hash table claims
a one-to-one slot assignment or a sorting argument matches inputs to outputs.
\end{focusbox}

\begin{example}[Injective but not surjective]
\label{ch38:ex-inj}
Let $f : \N \to \N$ be $f(n) = 2n$. Different inputs clearly give different outputs,
so $f$ is injective. But the image is $\{0,2,4,6,\dots\}$ --- the odd numbers are
never produced --- so $f$ is not surjective. A function can be one without the other.
\end{example}

The three properties are easiest to see as pictures of arrows from domain to
codomain. Injective forbids two arrows landing on the same target; surjective demands
every target be hit; a bijection is both at once --- a perfect pairing with no target
missed and none doubled.

\begin{center}
\begin{tikzpicture}[scale=1.0,every node/.style={font=\scriptsize},
    dot/.style={circle,fill,inner sep=1.6pt},
    >=Stealth]
  % --- Injective, not surjective ---
  \begin{scope}
    \node[font=\small\bfseries] at (0.5,2.0) {Injective};
    \foreach \i in {1,2,3} \node[dot,sourceblue] (L\i) at (0,{2.0-\i*0.55}) {};
    \foreach \i in {1,2,3,4} \node[dot,targetgreen] (R\i) at (1,{2.3-\i*0.55}) {};
    \draw[->,structurepurple] (L1) -- (R1);
    \draw[->,structurepurple] (L2) -- (R2);
    \draw[->,structurepurple] (L3) -- (R4);
  \end{scope}
  % --- Surjective, not injective ---
  \begin{scope}[xshift=3.6cm]
    \node[font=\small\bfseries] at (0.5,2.0) {Surjective};
    \foreach \i in {1,2,3,4} \node[dot,sourceblue] (S\i) at (0,{2.3-\i*0.55}) {};
    \foreach \i in {1,2,3} \node[dot,targetgreen] (T\i) at (1,{2.0-\i*0.55}) {};
    \draw[->,structurepurple] (S1) -- (T1);
    \draw[->,structurepurple] (S2) -- (T1);
    \draw[->,structurepurple] (S3) -- (T2);
    \draw[->,structurepurple] (S4) -- (T3);
  \end{scope}
  % --- Bijective ---
  \begin{scope}[xshift=7.2cm]
    \node[font=\small\bfseries] at (0.5,2.0) {Bijective};
    \foreach \i in {1,2,3} \node[dot,sourceblue] (U\i) at (0,{1.75-\i*0.55}) {};
    \foreach \i in {1,2,3} \node[dot,targetgreen] (V\i) at (1,{1.75-\i*0.55}) {};
    \draw[->,structurepurple] (U1) -- (V1);
    \draw[->,structurepurple] (U2) -- (V2);
    \draw[->,structurepurple] (U3) -- (V3);
  \end{scope}
\end{tikzpicture}
\end{center}
In the left picture a codomain dot is left untouched (not surjective); in the middle,
two arrows share a target (not injective); only the right picture is both.

\section{Counting: the art of the answer without enumeration}
\label{ch38:counting}

Counting is where set language pays off. The question is always the same shape: ``how
many objects of such-and-such kind are there?'' If we had to list them all and tally,
counting would be useless --- there are often astronomically many. The whole skill is
finding the answer \emph{without} listing.

Everything rests on two rules.

\begin{formalbox}
\textbf{Sum rule.} If a set $S$ is split into disjoint pieces
$S_1, S_2, \dots, S_k$ (no overlap, together they cover all of $S$), then
\[
\abs{S} = \abs{S_1} + \abs{S_2} + \cdots + \abs{S_k}.
\]

\textbf{Product rule.} If an object is built by making a sequence of $k$ independent
choices, where the $i$-th choice has $n_i$ options \emph{regardless} of the earlier
choices, then the number of objects is
\[
n_1 \cdot n_2 \cdots n_k.
\]
\end{formalbox}

The product rule is exactly the grid from Example~\ref{ch38:ex-product} stretched to
$k$ dimensions. Want to count $3$-letter strings over a $26$-letter alphabet? Three
choices, $26$ options each, independent: $26^3 = 17{,}576$. Want license plates of
the form letter-letter-digit-digit-digit? $26 \cdot 26 \cdot 10 \cdot 10 \cdot 10 =
676{,}000$.

\begin{labbox}
Trace the product rule by hand. Count the binary strings of length $3$. Choice $1$:
$\{0,1\}$. Choice $2$: $\{0,1\}$. Choice $3$: $\{0,1\}$. The product rule predicts
$2^3 = 8$. Now actually write them out: $000, 001, 010, 011, 100, 101, 110, 111$.
Eight, as promised. Notice the structure: the first half all start with $0$, the
second half with $1$ --- that is the sum rule ($4 + 4$) sitting inside the product
rule. The two rules are always working together.
\end{labbox}

\subsection{Permutations: counting orderings}

A \emph{permutation} is an arrangement of objects in a row, where order matters and
nothing repeats. How many ways can you line up $n$ distinct objects?

Reason with the product rule. The first slot: $n$ choices. The second slot: $n-1$
remaining objects. The third: $n-2$. And so on down to the last slot with $1$ choice.
Multiply:
\[
n \cdot (n-1) \cdot (n-2) \cdots 2 \cdot 1 = n!,
\]
read ``$n$ factorial.'' By convention $0! = 1$ (there is exactly one way to arrange
nothing: do nothing).

\begin{example}[Three books on a shelf]
\label{ch38:ex-perm}
Three distinct books $\{A, B, C\}$ arranged in a row: $3! = 6$ ways. List them to be
sure: $ABC, ACB, BAC, BCA, CAB, CBA$. Six. The first letter has $3$ choices, then $2$,
then the last is forced --- $3 \cdot 2 \cdot 1 = 6$.
\end{example}

Sometimes you only fill $k$ of the $n$ slots --- arrange $k$ objects chosen from $n$,
order mattering. This is the \emph{$k$-permutation}, and the same product-rule
argument (now stopping after $k$ factors) gives
\[
P(n, k) = n \cdot (n-1) \cdots (n-k+1) = \frac{n!}{(n-k)!}.
\]
For instance, the number of ways to award gold, silver, and bronze among $8$ runners
is $P(8,3) = 8 \cdot 7 \cdot 6 = 336$.

\begin{aestheticbox}
Why does an algorithms book care about $n!$? Because $n!$ is the number of possible
orderings of $n$ items, and \emph{sorting} is the act of finding the one correct
ordering. Any comparison-based sorting algorithm must, in effect, distinguish among
all $n!$ possibilities, and a decision tree that distinguishes $n!$ leaves needs depth
at least $\log_2(n!) = \Theta(n \log n)$. That single counting fact is the famous
$\Omega(n\log n)$ lower bound for comparison sorting. Counting orderings is not
trivia; it is the floor under every sort you will ever write.
\end{aestheticbox}

\subsection{Combinations: counting selections}

Now order does \emph{not} matter. How many ways can you choose a $k$-element subset
from $n$ distinct objects? This count is written $\binom{n}{k}$, read ``$n$ choose
$k$,'' and it is the most important number in this chapter.

Here is the clean derivation. First count \emph{ordered} selections of $k$ objects:
that is $P(n,k) = \tfrac{n!}{(n-k)!}$. But every unordered subset of size $k$ has been
counted once for each of its $k!$ orderings. So to get the unordered count we divide
out the overcounting:
\[
\binom{n}{k} = \frac{P(n,k)}{k!} = \frac{n!}{k!\,(n-k)!}.
\]
This ``count with order, then divide by the orderings'' move is the single most
useful technique in combinatorics. Learn to spot the overcounting and you can count
almost anything.

\begin{example}[Choosing a committee]
\label{ch38:ex-comb}
From $5$ people, choose a committee of $2$ (no roles, so order is irrelevant):
\[
\binom{5}{2} = \frac{5!}{2!\,3!} = \frac{120}{2 \cdot 6} = 10.
\]
Sanity check against the ordered count: there are $5 \cdot 4 = 20$ ways to pick a
first then a second person, but each committee $\{a,b\}$ was counted twice (as $ab$
and as $ba$), so $20 / 2! = 10$. The two routes agree.
\end{example}

\begin{trapbox}
The first reflex of every beginner is to use a permutation when a combination is
called for, or vice versa. The deciding question is one sentence: \emph{does
rearranging the chosen items give a different outcome?} Gold/silver/bronze --- yes,
order matters, use $P(n,k)$. A handful of lottery numbers, a poker hand, a subset of
array indices --- no, order is irrelevant, use $\binom{n}{k}$. Ask the question
before you reach for a formula.
\end{trapbox}

\subsection{Identities you should feel, not memorize}

A few facts about $\binom{n}{k}$ come up so often that you want them in your bones.

\textbf{Symmetry.} $\displaystyle \binom{n}{k} = \binom{n}{n-k}$. Choosing the $k$ you
keep is the same act as choosing the $n-k$ you leave behind --- a perfect bijection
between subsets and their complements. No algebra needed.

\textbf{Pascal's rule.} $\displaystyle \binom{n}{k} = \binom{n-1}{k-1} +
\binom{n-1}{k}$. Fix one special element, say the last. A $k$-subset either contains
it ($\binom{n-1}{k-1}$ ways to fill the rest) or does not ($\binom{n-1}{k}$ ways).
Those two cases are disjoint and cover everything --- the sum rule in action. This is
also the recurrence that builds Pascal's triangle, where each entry is the sum of the
two above it:

\begin{center}
\begin{tikzpicture}[scale=0.62,every node/.style={font=\small}]
  \foreach \r/\vals in {0/{1}, 1/{1,1}, 2/{1,2,1}, 3/{1,3,3,1}, 4/{1,4,6,4,1}} {
    \foreach \v [count=\c from 0] in \vals {
      \node at ({\c*1.4 - \r*0.7}, {-\r*1.0}) {\v};
    }
  }
  \draw[->,structurepurple,thick] (-0.7,-2.55) -- (-0.05,-3.5);
  \draw[->,structurepurple,thick] (0.7,-2.55) -- (0.05,-3.5);
  \node[structurepurple!70!black,font=\scriptsize] at (2.9,-3.0) {$3+3=6$};
\end{tikzpicture}
\end{center}

\textbf{Sum of a row.} $\displaystyle \sum_{k=0}^{n} \binom{n}{k} = 2^n$. The left
side counts subsets of every size; the right side counts \emph{all} subsets directly,
since each of the $n$ elements is independently in or out (product rule,
$2 \cdot 2 \cdots 2$). Two ways to count the same thing must agree.

\subsection{The binomial theorem}

All of this assembles into one beautiful formula. When you expand $(x + y)^n$, you
multiply out $n$ copies of $(x+y)$ and collect terms. Each term picks either $x$ or
$y$ from each factor; a term with $x^k y^{n-k}$ comes from choosing $x$ in exactly $k$
of the $n$ factors --- and there are $\binom{n}{k}$ ways to choose which. Hence

\begin{formalbox}
\textbf{Binomial theorem.}
\[
(x + y)^n = \sum_{k=0}^{n} \binom{n}{k}\, x^k\, y^{n-k}.
\]
The numbers $\binom{n}{k}$ are exactly the coefficients --- which is why they are
called \emph{binomial} coefficients.
\end{formalbox}

Setting $x = y = 1$ recovers $\sum_k \binom{n}{k} = 2^n$ instantly; setting $x = 1,\,
y = -1$ gives $\sum_k (-1)^{n-k}\binom{n}{k} = 0$, a handy alternating identity. One
formula, many corollaries.

\begin{walkbox}
On your next walk, prove this to yourself without algebra: why is the number of
\emph{lattice paths} from the bottom-left to the top-right corner of an $a \times b$
grid, moving only right or up, equal to $\binom{a+b}{a}$? Hint: every path is a
sequence of $a$ ``right'' steps and $b$ ``up'' steps; a path is fully determined by
\emph{which} of the $a+b$ moves are the rights. Counting paths becomes counting
subsets. This single idea --- encode a complicated object as a string, then count
strings --- reappears all over algorithm design.
\end{walkbox}

\section{Probability: reasoning about chance}
\label{ch38:probability}

Randomized algorithms flip coins. Quicksort picks a random pivot; hashing scatters
keys with a random function; sampling estimates a quantity from a few draws. To
analyze them we need a precise language for ``how likely'' and ``on average.'' That
language is probability, and --- happily --- it is built directly on the counting you
just learned.

\subsection{Sample spaces and events}

\begin{intuitionbox}
Picture every possible outcome of a random experiment laid out as points. That
collection is the \emph{sample space} $\Omega$. An \emph{event} is just a subset of
$\Omega$ --- a question you can ask whose answer is yes or no, like ``did the die show
an even number?'' Probability is then a way of assigning a weight to each event so the
weights behave sensibly.
\end{intuitionbox}

A \emph{probability} is a function $\Pr$ assigning to each event $A \subseteq \Omega$
a number $\Pr(A)$, obeying three axioms.

\begin{formalbox}
\textbf{Axioms of probability (Kolmogorov).}
\begin{enumerate}
\item $\Pr(A) \ge 0$ for every event $A$. (No negative chances.)
\item $\Pr(\Omega) = 1$. (\emph{Something} happens, with certainty.)
\item If $A$ and $B$ are disjoint ($A \cap B = \emptyset$), then
$\Pr(A \cup B) = \Pr(A) + \Pr(B)$. (Disjoint chances add.)
\end{enumerate}
\end{formalbox}

From these three, useful consequences fall out with a line or two of reasoning:
\[
\Pr(\overline{A}) = 1 - \Pr(A),
\qquad
\Pr(\emptyset) = 0,
\qquad
\Pr(A \cup B) = \Pr(A) + \Pr(B) - \Pr(A \cap B).
\]
That last one --- \emph{inclusion--exclusion} --- corrects for the overlap you would
otherwise double-count, exactly mirroring the Venn diagram from
Section~\ref{ch38:sets}.

When all $\abs{\Omega}$ outcomes are \emph{equally likely} (a fair die, a uniformly
random pivot), probability reduces to counting:
\[
\Pr(A) = \frac{\abs{A}}{\abs{\Omega}} = \frac{\text{favorable outcomes}}{\text{total outcomes}}.
\]
This is why the counting section came first: in the equally-likely world,
\emph{probability is just a ratio of two counts}.

\begin{example}[A pair of dice]
\label{ch38:ex-dice}
Roll two fair six-sided dice. The sample space is the $36$ ordered pairs
$(i, j)$ with $1 \le i, j \le 6$, all equally likely. What is $\Pr(\text{sum} = 7)$?
Count the favorable pairs: $(1,6),(2,5),(3,4),(4,3),(5,2),(6,1)$ --- six of them. So
\[
\Pr(\text{sum} = 7) = \frac{6}{36} = \frac{1}{6}.
\]
And $\Pr(\text{sum} = 12)$? Only $(6,6)$ works, so $\tfrac{1}{36}$. The counting did
all the work; the probability is a one-line division.

\begin{center}
\begin{tikzpicture}[scale=0.52,every node/.style={font=\scriptsize}]
  % grid of 36 cells, first die i = row (1 bottom .. 6 top), second die j = column
  \foreach \i in {1,...,6} {
    \foreach \j in {1,...,6} {
      \pgfmathtruncatemacro{\s}{\i+\j}
      \ifnum\s=7
        \node[draw,fill=trapred!30,minimum size=0.85cm,inner sep=0pt]
          at (\j,\i) {\i,\j};
      \else
        \node[draw,fill=gray!6,minimum size=0.85cm,inner sep=0pt]
          at (\j,\i) {\i,\j};
      \fi
    }
  }
  % axis labels
  \node[targetgreen!60!black] at (3.5,7.1) {second die $j$};
  \node[sourceblue!60!black,rotate=90] at (-0.1,3.5) {first die $i$};
\end{tikzpicture}
\end{center}
Each of the $36$ cells is one equally-likely outcome $(i,j)$; the six shaded cells on
the anti-diagonal are exactly the rolls with $i + j = 7$, giving $6/36 = 1/6$.
\end{example}

\subsection{Conditional probability}

New information changes the odds. Once you learn that an event $B$ has occurred, the
sample space effectively shrinks to $B$, and you re-measure everything inside it.

\begin{formalbox}
The \emph{conditional probability} of $A$ given $B$ (with $\Pr(B) > 0$) is
\[
\Pr(A \mid B) = \frac{\Pr(A \cap B)}{\Pr(B)}.
\]
\end{formalbox}

The picture is clean: you are restricting attention to the world where $B$ happened
(divide by $\Pr(B)$ to renormalize so that world has total weight $1$), and asking how
much of that world also has $A$ (the numerator $\Pr(A \cap B)$).

\begin{example}[Dice, with a hint]
\label{ch38:ex-cond}
Roll two fair dice. Let $A$ be ``the sum is $7$'' and $B$ be ``the first die is $3$.''
We have $\Pr(B) = \tfrac{6}{36} = \tfrac16$. The pairs in $A \cap B$: only $(3,4)$
gives both a leading $3$ and a sum of $7$, so $\Pr(A \cap B) = \tfrac{1}{36}$. Thus
\[
\Pr(A \mid B) = \frac{1/36}{1/6} = \frac{1}{6}.
\]
Knowing the first die showed $3$ left the chance of a $7$ unchanged at $\tfrac16$ ---
a hint of the independence idea coming next.
\end{example}

Rearranging the definition gives the \emph{chain rule}
$\Pr(A \cap B) = \Pr(A \mid B)\,\Pr(B)$, and summing conditional pieces over a
partition $B_1, \dots, B_k$ of $\Omega$ gives the \emph{law of total probability}
\[
\Pr(A) = \sum_{i=1}^{k} \Pr(A \mid B_i)\,\Pr(B_i),
\]
which is how you compute a probability by splitting into cases --- the sum rule
wearing a probabilistic coat. Flipping the conditioning around yields \emph{Bayes'
rule}, $\Pr(B \mid A) = \Pr(A \mid B)\,\Pr(B)/\Pr(A)$, the engine behind every
``update my belief given evidence'' calculation.

\subsection{Independence}

\begin{intuitionbox}
Two events are \emph{independent} when learning one tells you nothing about the
other. Knowing the coin came up heads says nothing about tomorrow's weather; the two
live in separate worlds.
\end{intuitionbox}

Formally, $A$ and $B$ are \emph{independent} when
\[
\Pr(A \cap B) = \Pr(A)\,\Pr(B),
\]
equivalently $\Pr(A \mid B) = \Pr(A)$ whenever $\Pr(B) > 0$ (the condition did not move
the needle). In Example~\ref{ch38:ex-cond} the sum-is-$7$ event was independent of the
first die's value --- a small surprise, but the formula confirms it.

\begin{trapbox}
Independence is a property of the \emph{numbers}, not a feeling. ``These seem
unrelated'' is not a proof. You must check $\Pr(A \cap B) = \Pr(A)\Pr(B)$. And beware:
events can be \emph{pairwise} independent yet not \emph{mutually} independent --- the
three-way relationship can fail even when every pair checks out. When an analysis
needs full independence (as in some hashing arguments), pairwise is not enough, and
sloppily assuming it has sunk many a proof.
\end{trapbox}

\section{Random variables and expectation}
\label{ch38:rv}

\subsection{What a random variable is}

\begin{intuitionbox}
A \emph{random variable} is a number that depends on the random outcome --- a
measurement you read off after the experiment runs. ``The number of comparisons
quicksort makes'' is a random variable: it has no single fixed value, because it
depends on which random pivots came up.
\end{intuitionbox}

Formally a (discrete) random variable $X$ is a function $X : \Omega \to \R$. For each
value $v$ it can take, the probability $\Pr(X = v)$ is shorthand for the probability of
the event $\{\, \omega \in \Omega : X(\omega) = v \,\}$. The list of these
probabilities, over all values, is the \emph{distribution} of $X$.

\subsection{Expectation: the average value}

The \emph{expected value} (or mean) of $X$ is the average of its values, each weighted
by how likely it is:
\[
\E[X] = \sum_{v} v \cdot \Pr(X = v),
\]
the sum running over all values $X$ can take. Think of it as the center of mass of the
distribution --- the balance point if you placed weight $\Pr(X=v)$ at position $v$ on a
ruler.

\begin{example}[Expected value of one die]
\label{ch38:ex-die-exp}
For a fair die, $X \in \{1,2,3,4,5,6\}$ each with probability $\tfrac16$:
\[
\E[X] = \frac{1}{6}(1 + 2 + 3 + 4 + 5 + 6) = \frac{21}{6} = 3.5.
\]
Notice the average, $3.5$, is a value the die can never actually show. That is fine:
expectation is a long-run average, not a prediction of any single roll.
\end{example}

\subsection{Linearity of expectation: the workhorse}

Here is the single most useful fact in all of randomized algorithm analysis.

\begin{formalbox}
\textbf{Linearity of expectation.} For \emph{any} random variables $X$ and $Y$ and
constants $a, b$,
\[
\E[aX + bY] = a\,\E[X] + b\,\E[Y].
\]
This holds \emph{whether or not $X$ and $Y$ are independent.} It extends to any finite
sum: $\E[X_1 + \cdots + X_n] = \E[X_1] + \cdots + \E[X_n]$.
\end{formalbox}

The ``no independence required'' clause is what makes linearity magical. You can break
a horrendously tangled total into simple pieces, take each piece's expectation in
isolation, and add --- never once worrying about how the pieces interact.

\begin{focusbox}
The standard recipe: to find $\E[X]$ where $X$ counts how many times something
happens, write $X = X_1 + X_2 + \cdots + X_n$ where each $X_i$ is an \emph{indicator}
--- a $0/1$ random variable that is $1$ exactly when the $i$-th thing happens. The key
identity is
\[
\E[X_i] = 0 \cdot \Pr(X_i = 0) + 1 \cdot \Pr(X_i = 1) = \Pr(X_i = 1).
\]
\emph{The expectation of an indicator is just the probability of the event it
indicates.} So by linearity, $\E[X] = \sum_i \Pr(\text{$i$-th event happens})$. A hard
average becomes a sum of easy probabilities. This is \emph{the} technique behind
quicksort's expected-comparison count, hashing's expected chain length, and dozens
more.
\end{focusbox}

\begin{example}[Expected number of fixed points]
\label{ch38:ex-fixedpoints}
Shuffle $n$ cards uniformly at random. How many cards land in their original position
on average? Let $X$ be that count, and let $X_i = 1$ if card $i$ stays put. Card $i$ is
equally likely to land in any of the $n$ positions, so $\Pr(X_i = 1) = \tfrac1n$, hence
$\E[X_i] = \tfrac1n$. The cards' fates are tangled (if one is fixed, the others are
constrained), but linearity does not care:
\[
\E[X] = \sum_{i=1}^{n} \E[X_i] = \sum_{i=1}^{n} \frac1n = n \cdot \frac1n = 1.
\]
On average, exactly \emph{one} card stays home, no matter how big the deck. We never
touched the distribution of $X$ --- only $n$ trivial indicators. That is the power of
linearity.
\end{example}

\begin{labbox}
Verify Example~\ref{ch38:ex-fixedpoints} by brute force for $n = 3$. The $3! = 6$
permutations of $(1,2,3)$, with their fixed-point counts:
\[
\begin{array}{llll}
123 \to 3, & 132 \to 1, & 213 \to 1, \\
231 \to 0, & 312 \to 0, & 321 \to 1.
\end{array}
\]
Average $= (3 + 1 + 1 + 0 + 0 + 1)/6 = 6/6 = 1$. Exactly the predicted answer, and you
can see why no shortcut would have been faster than linearity once $n$ grows.
\end{labbox}

The indicator recipe is worth one picture. The tangled total $X$ splits into a row of
simple $0/1$ flags, each easy to average on its own; linearity lets us sum those
averages even though the flags are correlated.

\begin{center}
\begin{tikzpicture}[scale=1.0,every node/.style={font=\scriptsize}]
  \node[font=\small] at (-1.2,0) {$X \;=\;$};
  \foreach \i [count=\c from 0] in {1,2,3,4} {
    \node[draw,sourceblue!60!black,fill=sourceblue!8,rounded corners,
          minimum width=0.9cm,minimum height=0.7cm] (b\i) at (\c*1.5,0) {$X_{\i}$};
    \ifnum\i<4 \node at (\c*1.5+0.75,0) {$+$}; \fi
  }
  \foreach \i in {1,2,3,4} {
    \node[font=\scriptsize,gray] at (b\i.south) [below=2pt]
      {$\E[X_{\i}]=\Pr(\text{event }\i)$};
  }
  \node[font=\scriptsize] at (6.1,0.0) {$+\cdots$};
\end{tikzpicture}
\end{center}

\subsection{Variance, briefly}

Expectation tells you the center; \emph{variance} tells you the spread. It is the
expected squared distance from the mean,
\[
\Var[X] = \E\!\big[(X - \E[X])^2\big] = \E[X^2] - (\E[X])^2,
\]
and its square root $\sqrt{\Var[X]}$ is the \emph{standard deviation}. Unlike
expectation, variance is \emph{not} linear in general --- but it does add for
\emph{independent} variables: $\Var[X+Y] = \Var[X] + \Var[Y]$ when $X, Y$ are
independent. Variance is what you reach for when you want to argue that a randomized
algorithm is not just good \emph{on average} but good \emph{almost always}.

\section{Useful inequalities}
\label{ch38:inequalities}

Often you cannot compute a probability exactly --- but you can \emph{bound} it, and a
bound is usually all an algorithm analysis needs (``the failure probability is at most
$1/n^2$'' is plenty). Three inequalities do most of the heavy lifting.

\subsection{Markov's inequality}

\begin{formalbox}
\textbf{Markov's inequality.} If $X \ge 0$ and $t > 0$, then
\[
\Pr(X \ge t) \le \frac{\E[X]}{t}.
\]
\end{formalbox}

The intuition: a nonnegative quantity cannot often be far above its average, because
those rare large values would, by themselves, drag the average up past where it
actually sits. Equivalently, taking $t = c\,\E[X]$: a nonnegative random variable is at
least $c$ times its mean with probability at most $1/c$. It needs nothing but
nonnegativity and the mean --- which is exactly why it is so widely applicable, and so
loose.

\begin{example}[A quick Markov bound]
\label{ch38:ex-markov}
An algorithm runs in expected time $\E[T] = 100$ steps. What can we say about the
chance it takes $1000$ or more? With $t = 1000$,
\[
\Pr(T \ge 1000) \le \frac{100}{1000} = \frac{1}{10}.
\]
So at most a $10\%$ chance of being ten times over budget --- derived from the mean
alone, knowing nothing else about $T$.
\end{example}

\begin{center}
\begin{tikzpicture}[scale=1.0,every node/.style={font=\scriptsize},>=Stealth]
  \draw[->] (-0.3,0) -- (8.3,0) node[right,font=\small] {$X$};
  \draw[thick] (0,-0.08) -- (0,0.08) node[above=1pt] {$0$};
  % mean marker
  \draw[thick,sourceblue] (2,-0.12) -- (2,0.12);
  \node[sourceblue!70!black,above=2pt] at (2,0.12) {$\E[X]$};
  % threshold t
  \draw[thick,trapred] (6,-0.12) -- (6,0.12);
  \node[trapred,above=2pt] at (6,0.12) {$t$};
  % shaded tail region
  \fill[trapred!18] (6,-0.55) rectangle (8.0,-0.05);
  \draw[trapred] (6,-0.55) rectangle (8.0,-0.05);
  \node[trapred,font=\scriptsize] at (7.0,-0.3) {$X \ge t$};
  \node[align=center,font=\scriptsize] at (7.0,-1.05)
    {total weight here\\$\le \E[X]/t$};
\end{tikzpicture}
\end{center}
Markov says the probability mass to the right of $t$ is capped: a nonnegative variable
with mean $\E[X]$ cannot pile too much weight far above its average without dragging
the average up past where it sits.

\subsection{Chebyshev's inequality}

When you also know the variance, you can bound deviations on \emph{both} sides of the
mean, and tightly.

\begin{formalbox}
\textbf{Chebyshev's inequality.} For any random variable $X$ with finite variance, and
any $t > 0$,
\[
\Pr\!\big(\abs{X - \E[X]} \ge t\big) \le \frac{\Var[X]}{t^2}.
\]
\end{formalbox}

It is just Markov applied to the nonnegative variable $(X - \E[X])^2$. The payoff: the
probability of straying $k$ standard deviations from the mean is at most $1/k^2$ ---
small spread means $X$ clusters tightly around $\E[X]$, which is precisely the
``good almost always'' guarantee variance was promising.

\subsection{The union bound}

\begin{formalbox}
\textbf{Union bound (Boole's inequality).} For any events $A_1, \dots, A_n$,
\[
\Pr\!\Big( \bigcup_{i=1}^{n} A_i \Big) \le \sum_{i=1}^{n} \Pr(A_i).
\]
\end{formalbox}

Dirt simple, endlessly useful. If each of $n$ ``bad'' events is individually unlikely,
the chance that \emph{any} of them happens is at most the sum of their probabilities.
It is the inclusion--exclusion identity with the (subtracted) overlaps thrown away ---
which can only make the right side bigger, so it is always a valid upper bound.

\begin{example}[Union bound in action]
\label{ch38:ex-union}
Suppose a randomized data structure has $n$ operations, and each individually fails
with probability at most $1/n^2$. Then the probability that \emph{some} operation
fails is at most
\[
\sum_{i=1}^{n} \frac{1}{n^2} = n \cdot \frac{1}{n^2} = \frac{1}{n}.
\]
So with probability at least $1 - 1/n$ --- ``with high probability'' --- the entire run
is flawless. This pattern, bound each rare event then union them, is how essentially
every high-probability guarantee in the book gets proved.
\end{example}

\begin{historybox}
This material lines up with Appendices B and C of CLRS. \textbf{Appendix B} covers
sets, relations, functions, and the basics of counting; our
Sections~\ref{ch38:sets}--\ref{ch38:relations} and the start of
Section~\ref{ch38:counting} are the gentle on-ramp to it. \textbf{Appendix C} covers
counting (permutations and combinations), discrete probability, random variables,
expectation, and the standard inequalities; our
Sections~\ref{ch38:counting}--\ref{ch38:inequalities} mirror it section for section.
The historical roots run deep: the systematic study of counting and chance began with
the seventeenth-century correspondence between Blaise Pascal and Pierre de Fermat over
a gambling problem, and the modern axioms in Section~\ref{ch38:probability} are due to
Andrey Kolmogorov in 1933. When you reach Chapter 5 (probabilistic analysis and
randomized algorithms), Chapter 7 (quicksort), and Chapter 11 (hashing) in CLRS, you
will be cashing in everything assembled here.
\end{historybox}

\begin{complexitybox}
The counting and probability facts you will reuse most:
\begin{center}
\begin{tabular}{@{}ll@{}}
\toprule
\textbf{Quantity} & \textbf{Formula} \\
\midrule
Subsets of an $n$-set & $2^n$ \\
Orderings of $n$ items & $n!$ \\
Ordered $k$ from $n$ & $P(n,k) = n!/(n-k)!$ \\
Unordered $k$ from $n$ & $\binom{n}{k} = n!\,/\,(k!\,(n-k)!)$ \\
Equally-likely probability & $\Pr(A) = \abs{A}/\abs{\Omega}$ \\
Conditional probability & $\Pr(A\mid B) = \Pr(A\cap B)/\Pr(B)$ \\
Expectation & $\E[X] = \sum_v v\,\Pr(X=v)$ \\
Linearity & $\E[X+Y] = \E[X] + \E[Y]$ (always) \\
Indicator expectation & $\E[X_i] = \Pr(\text{event } i)$ \\
Markov & $\Pr(X \ge t) \le \E[X]/t$ \\
Chebyshev & $\Pr(\abs{X-\E[X]}\ge t) \le \Var[X]/t^2$ \\
Union bound & $\Pr(\bigcup A_i) \le \sum \Pr(A_i)$ \\
\bottomrule
\end{tabular}
\end{center}
\end{complexitybox}

\section{Exercises}
\label{ch38:exercises}

\begin{exercise}[Warmup: set arithmetic]
Let $A = \{1,2,3,4,5\}$ and $B = \{2,4,6\}$ with universe $U = \{1,2,\dots,8\}$.
Compute $A \cup B$, $A \cap B$, $A \setminus B$, $\overline{A}$, and $\abs{A \times
B}$. Then verify De Morgan's law $\overline{A \cup B} = \overline{A} \cap \overline{B}$
by computing both sides.
\end{exercise}

\begin{exercise}[Type-checking $\in$ versus $\subseteq$]
For the set $S = \{1, \{2\}, \{1,2\}\}$, decide whether each statement is true or
false, and explain in one sentence: (a) $1 \in S$; (b) $\{2\} \in S$; (c)
$2 \in S$; (d) $\{1\} \subseteq S$; (e) $\{\{2\}\} \subseteq S$.
\end{exercise}

\begin{exercise}[Permutation versus combination]
A standard deck has $52$ cards. (a) How many ways are there to deal a $5$-card poker
hand, where order does not matter? (b) How many ways to deal $5$ cards in order to five
seated players, one each, where who gets which card matters? Express each as a
product/factorial expression and as a number, and state which is larger and why.
\end{exercise}

\begin{exercise}[Counting strings]
How many strings of length $4$ over the alphabet $\{a,b,c\}$ contain at least one
$a$? (Hint: it is far easier to count the complement --- strings with \emph{no} $a$ ---
and subtract from the total. This complement trick is worth internalizing.)
\end{exercise}

\begin{exercise}[Pascal by hand]
Use Pascal's rule $\binom{n}{k} = \binom{n-1}{k-1} + \binom{n-1}{k}$ to build the row
of $\binom{6}{k}$ for $k = 0, 1, \dots, 6$, starting from the row for $n = 5$
($1,5,10,10,5,1$). Then check that your row sums to $2^6 = 64$.
\end{exercise}

\begin{exercise}[Conditional dice]
Roll two fair dice. Given that the sum is $8$, what is the probability that at least one
die shows a $5$? Identify the conditioning event $B$, list the outcomes in $B$, then
list those that also satisfy $A$, and compute $\Pr(A \mid B)$ directly.
\end{exercise}

\begin{exercise}[Independence check]
Roll two fair dice. Let $A$ be ``the first die is even'' and $B$ be ``the sum is
$7$.'' Compute $\Pr(A)$, $\Pr(B)$, and $\Pr(A \cap B)$, and determine whether $A$ and
$B$ are independent. Does the answer match your intuition?
\end{exercise}

\begin{exercise}[Linearity to the rescue]
You roll a fair die $n$ times. Let $X$ be the number of times a $6$ appears. Using
indicator variables $X_i$ (one per roll) and linearity of expectation, find $\E[X]$.
Then, separately, find the expected number of times a $6$ is \emph{immediately
followed by another $6$} among the $n-1$ adjacent pairs.
\end{exercise}

\begin{exercise}[Expected maximum of two dice]
Roll two fair dice and let $M$ be the larger of the two values (the maximum). Write out
the distribution of $M$ (the probability $\Pr(M = v)$ for $v = 1, \dots, 6$) and
compute $\E[M]$. (Hint: $\Pr(M \le v) = (v/6)^2$, since both dice must be at most $v$;
turn that into $\Pr(M = v)$.)
\end{exercise}

\begin{exercise}[Choosing the right inequality]
A randomized algorithm has expected running time $\E[T] = 50$ and variance
$\Var[T] = 25$. (a) Use Markov's inequality to bound $\Pr(T \ge 200)$. (b) Use
Chebyshev's inequality to bound $\Pr(\abs{T - 50} \ge 30)$. (c) In one or two
sentences, explain why Chebyshev gave a more informative bound here, and what extra
information made that possible.
\end{exercise}

\begin{exercise}[Union bound design]
\label{ch38:ex-uniondesign}
You hash $n$ keys into a table. For a fixed pair of distinct keys, suppose the
probability they collide (land in the same slot) is exactly $1/m$, where $m$ is the
table size. Using the union bound over all pairs, show that the probability of
\emph{any} collision is at most $\binom{n}{2}/m$. Then determine how large $m$ must be,
as a function of $n$, to make this bound at most $1/2$. (This is the birthday-paradox
threshold, and it explains why hash tables need room to breathe.)
\end{exercise}

\begin{exercise}[Challenge: a counting bijection]
Prove that the number of ways to write a positive integer $n$ as an ordered sum of
positive integers (a \emph{composition}, e.g. $4 = 1+3 = 3+1 = 2+2 = 1+1+2 = \cdots$)
is exactly $2^{n-1}$. (Hint: line up $n$ dots in a row; between each adjacent pair of
dots there is a gap, and you either ``cut'' there or not. Set up a bijection between
compositions and subsets of the $n-1$ gaps, then count the subsets.)
\end{exercise}
